\documentclass[../../main.tex]{subfiles}% 注意这里的文件路径不能用 ./main.tex ,否则用latexmk编译子文件会报错
\graphicspath{{\subfix{./image/}}{\subfix{../../image/}}} % 指定图片引用路径,后续可以直接使用图片文件名
% 如果图片存放在上级目录,则使用 ../../image/ 作为备用路径

% 例如:
% \begin{figure}[H]
% \centering
% \includegraphics[width=0.3\textwidth]{图.png}
% \caption{}
% \label{figure:图}
% \end{figure}
% 注意:上述\label{}一定要放在\caption{}之后,否则引用图片序号会只会显示??.

\begin{document}

\section{Riemann映射定理}

\begin{theorem}[Riemann映射定理]\label{theorem:Riemann映射定理}
设 \(G\) 是 \(\mathbb{C}\) 中的单连通域, \(G \neq \mathbb{C}\). 对于 \(G\) 中的任意点 \(a\), 存在唯一的函数 \(f: G \to \mathbb{C}\), 使得
\begin{enumerate}[(1)]
\item \(f\) 在 \(G\) 中全纯且单叶;
\item \(f(a) = 0\), \(f'(a) > 0\);
\item \(f(G) = B(0, 1)\).
\end{enumerate}
\end{theorem}
\begin{proof}
\begin{enumerate}[(1)]
\item 作一函数 \(\varphi: G \to \mathbb{C}\), 它在 \(G\) 中全纯且单叶, 使得 \(\varphi(G)\) 是一有界域. 事实上, 因为 \(G \neq \mathbb{C}\), 故存在 \(b \in G, b \neq \infty\). 在 \(G\) 中取 \(\sqrt{z-b}\) 的单值全纯分支, 记为 \(g(z) = \sqrt{z-b}\). 显然, 它在 \(G\) 中是单叶的. 设 \(g(G) = E\), 则由\hyperref[theorem:定理4.4.6-开映射定理]{开映射定理}知, \(E\) 是 \(\mathbb{C}\) 中的域. \(E\) 有一个简单的性质: 如果 \(w \in E\), 那么 \(-w \notin E\). 因为如果 \(-w \in E\), 则有 \(z_1, z_2 \in G\), 使得 \(\sqrt{z_1-b} = w, \sqrt{z_2-b} = -w\). 两边平方后即得 \(z_1 = z_2\), 因而 \(w = -w\), 即 \(w = 0\). 这是不可能的, 因为 \(g\) 在 \(G\) 中没有零点. 因为 \(a \in G\), 所以 \(g(a)\) 是 \(E\) 的内点, 故存在 \(\delta > 0\), 使得 \(B(g(a), \delta) \subseteq E\), 因而 \(B(-g(a), \delta) \subseteq \mathbb{C} \setminus E\). 于是, 当 \(z \in G\) 时, 便有 \(|g(z) + g(a)| > \delta\). 令
\begin{align*}
\varphi(z) = \frac{1}{g(z) + g(a)},
\end{align*}
则当 \(z \in G\) 时, \(|\varphi(z)| < \frac{1}{\delta}\), 即 \(\varphi(G) \subset B\left(0, \frac{1}{\delta}\right)\), 这说明 \(\varphi(G)\) 是一个有界域.

\item 由步骤(1)所证, 我们不妨假定 \(G\) 是有界域. 在 \(G\) 上定义函数族
\begin{align*}
\mathscr{F} = \{ f : \text{\(f\) 在 \(G\) 上全纯且单叶}, f(a) = 0, f'(a) > 0, f(G) \subseteq B(0, 1) \}.
\end{align*}
我们首先证明 \(\mathscr{F}\) 不是空集. 事实上, 因为 \(G\) 是有界域, 故存在 \(R > 0\), 使得 \(G \subset B(0, R)\). 若令
\begin{align*}
f(z) = \frac{1}{2R} (z - a),
\end{align*}
由于 \(|f(z)| \leqslant \frac{1}{2R} (|z| + |a|) < 1\), 故 \(f \in \mathscr{F}\), 这证明 \(\mathscr{F}\) 非空.
现取 \(r > 0\), 使得 \(B(a, r) \subset G\). 根据 Cauchy 不等式, 对任意 \(f \in \mathscr{F}\), 有
\begin{align*}
f'(a) < \frac{1}{r} \sup\{ |f(z)| : z \in B(a, r) \} < \frac{1}{r}.
\end{align*}
这说明数集 \(\{ f'(a) : f \in \mathscr{F} \}\) 有上界, 设 \(M\) 为其上确界, 则
\begin{align*}
M = \sup\{ f'(a) : f \in \mathscr{F} \} < \infty.
\end{align*}
下面我们证明, 存在 \(f_* \in \mathscr{F}\), 使得 \(f_*'(a) = M\). 事实上, 存在 \(f_n \in \mathscr{F}\), 使得 \(\lim\limits_{n \to \infty} f_n'(a) = M\). 由于 \(\mathscr{F}\) 在 \(G\) 上一致有界, 因而由\hyperref[theorem:复变-----Montel定理]{Montel定理}, 它是一个正规族, 故在 \(\{ f_n \}\) 中可以取出在 \(G\) 上内闭一致收敛的子列 \(\{ f_{n_k} \}\), 设其极限函数为 \(f_*\), 即
\begin{align}
\lim\limits_{k \to \infty} f_{n_k}(z) = f_*(z). \label{eq:1a2b3c}
\end{align}
由\cref{corollary:复变函数--函数项级数Weierstrass定理推论}, \(f_* \in H(G)\), 且
\begin{align}
\lim\limits_{k \to \infty} f'_{n_k}(z) = f'_*(z). \label{eq:4d5e6f}
\end{align}
由 \eqref{eq:4d5e6f} 及 \(\lim\limits_{k \to \infty} f'_{n_k}(a) = M\), 即得
\begin{align}
f'_*(a) = M. \label{eq:7g8h9i}
\end{align}
现在证明 \(f_* \in \mathscr{F}\). 因为每一个 \(f_{n_k}\) 都是 \(G\) 中的单叶全纯函数, 由 \eqref{eq:7g8h9i} 式知 \(f'_*\) 不是常数, 由\cref{theorem:定理4.4.11}知道 \(f_*\) 也是 \(G\) 中的单叶全纯函数. 由 \eqref{eq:1a2b3c} 式易知 \(f_*(a) = 0\), 且 \(|f_*(z)| \leqslant 1, z \in G\). 由\hyperref[theorem:最大模原理-定理4.5.1]{最大模原理}, \(|f_*(z)| = 1\) 不能成立, 因而只有 \(|f_*(z)| < 1\). 这就证明了 \(f_* \in \mathscr{F}\).

\item 现在证明 \(f_*\) 即是定理中要找的函数, 即要证明 \(f_*(G) = B(0, 1)\). 如果不是这样, 设
\begin{align*}
f_*(G) = G_1 \subset B(0, 1),
\end{align*}
那么 \(G_1\) 是 \(B(0, 1)\) 中的一个单连通域, 但 \(G_1 \neq B(0, 1)\). 因而存在 \(u_0 \in B(0, 1), u_0 \notin G_1\), 由此即可引出一个矛盾. 为此, 令
\begin{align*}
v = \varphi_{u_0}(u) \triangleq \frac{u_0 - u}{1 - \overline{u_0} u},
\end{align*}
于是
\begin{align*}
\varphi_{u_0}(u_0) = 0, \quad
\varphi'_{u_0}(0) = - (1 - |u_0|^2).
\end{align*}
\(\varphi_{u_0}\) 把 \(G_1\) 映为 \(B(0, 1)\) 中的单连通域 \(G_2\), 把 \(u_0\) 映为原点, 把原点映为 \(v_0 = u_0\), 所以 \(0 \notin G_2\). 再令
\begin{align*}
s = p(v) \triangleq \sqrt{v},
\end{align*}
于是
\begin{align*}
s_0 = p(v_0) \triangleq \sqrt{v_0}, \quad
p'(v_0) = \frac{1}{2 \sqrt{v_0}}.
\end{align*}
它在 \(G_2\) 中能分出单值全纯的分支, 记 \(p(G_2) = G_3\), 则 \(G_3\) 是 \(B(0, 1)\) 中的单连通域, \(0 \notin G_3\), 最后, 令
\begin{align*}
w = q(s) \triangleq \frac{s_0}{|s_0|} \varphi_0(s),
\end{align*}
于是
\begin{align*}
q'(s_0) = - \frac{s_0}{|s_0|} \frac{1}{1 - |s_0|^2}.
\end{align*}
它把 \(G_3\) 映为 \(B(0, 1)\) 中的单连通域 \(G_4\), 把 \(s_0\) 映为原点. 于是, 复合函数
\begin{align*}
w = (q \circ p \circ \varphi_{u_0} \circ f_*)(z) \triangleq w(z)
\end{align*}
把 \(G\) 映为 \(G_4\), 且 \(w(a) = 0\) (见\cref{figure:复变图7.6}). \(w(z)\) 显然是 \(G\) 上的单叶全纯函数. 现在来计算 \(w'(a)\), 根据复合函数的求导法则, 有
\begin{align*}
w'(a) &= q'(s_0) p'(v_0) \varphi'_{u_0}(0) f'_*(a) \\
&= \frac{s_0}{|s_0|} \frac{1}{1 - |s_0|^2} \frac{1}{2 \sqrt{v_0}} (1 - |u_0|^2) M \\
&= \frac{1 - |u_0|^2}{2 |s_0| (1 - |s_0|^2)} M > 0.
\end{align*}
所以 \(w \in \mathscr{F}\). 若记 \(u_0 = \rho e^{i\theta}\), 则 \(v_0 = - u_0 = \rho e^{i(\theta + \pi)}\), \(|s_0| = |\sqrt{v_0}| = \sqrt{\rho}\), \(1 - |s_0|^2 = 1 - \rho\), 因而
\begin{align*}
w'(a) = \frac{1 - \rho^2}{2 \sqrt{\rho} (1 - \rho)} M = \frac{1 + \rho}{2 \sqrt{\rho}} M > M.
\end{align*}
\begin{figure}[H]
\centering
\includegraphics[width=0.8\textwidth]{复变图7.6.png}
\caption{}
\label{figure:复变图7.6}
\end{figure}
这就是说, 我们在 \(\mathscr{F}\) 中找到了一个函数 \(w\), 它在 \(a\) 点的导数大于 \(M\), 这是不可能的. 这个矛盾说明假设不成立, 于是 \(f_*(G) = B(0, 1)\).

\item 最后证明满足定理中三个条件的 \(f\) 是唯一的. 如果还有函数 \(g\) 也满足这三个条件, 令
\begin{align*}
h = f \circ g^{-1},
\end{align*}
那么 \(h(B(0, 1)) = B(0, 1)\), 因而 \(h \in \operatorname{Aut}(B(0, 1))\). 又因 \(h(0) = 0\), 因而由\cref{theorem:定理4.5.5}知道
\begin{align*}
h(z) = e^{i\theta} z.
\end{align*}
但由于 \(h'(0) = \frac{f'(a)}{g'(a)} > 0\), 即知 \(h(z) = z\), 从而
\begin{align*}
f(z) \equiv g(z), \quad z \in G.
\end{align*}
\end{enumerate}
\end{proof}

\begin{corollary}
设 \(G\) 和 \(D\) 是 \(\mathbb{C}\) 中的两个单连通域, 如果它们都不是整个平面 \(\mathbb{C}\), 那么对于给定的 \(z_0 \in G\) 和 \(w_0 \in D\), 存在唯一的函数 \(f\), 它在 \(G\) 上单叶且全纯, \(f(z_0) = w_0\), \(f'(z_0) > 0\), 且 \(f(G) = D\).
\end{corollary}
\begin{remark}
这个定理断言, 两个异于 \(\mathbb{C}\) 的单连通域是全纯等价的. 

已经知道单连通域和多连通域一定不全纯等价, 那么两个 \(n\) (\(n > 1\)) 连通域是不是全纯等价呢? 一般来说, 答案是否定的. \cref{proposition:例6.1.6复变}已为我们提供了这样的例子: 即使两个最简单的二连通域——两个同心圆环 \(D = \{ z \in \mathbb{C} : r_1 < |z| < r_2 \}\) 和 \(G = \{ w \in \mathbb{C} : R_1 < |w| < R_2 \}\) 全纯等价也是有条件的, 要求它们满足 \(\frac{r_2}{r_1} = \frac{R_2}{R_1}\). 这从另一个角度说明了 Riemann 映射定理的重要性.
\end{remark}
\begin{proof}
由\hyperref[theorem:Riemann映射定理]{Riemann映射定理}, 存在唯一的全纯函数 \(g\), 使得 \(\zeta = g(z)\) 把 \(G\) 一一地映为 \(|\zeta| < 1\), 且 \(g(z_0) = 0, g'(z_0) > 0\). 同样道理, 存在唯一的全纯函数 \(\zeta = h(w)\), 它把 \(D\) 一一地映为 \(|\zeta| < 1\), 且 \(h(w_0) = 0, h'(w_0) > 0\). 于是, 函数 \(w = (h^{-1} \circ g)(z) \triangleq f(z)\) 即是符合定理条件的函数.
\end{proof}



\end{document}